\section{Discussion}
\label{sec:discussion}
What follows applies to the archive of a single publisher, read in an exploratory way (Section~\ref{sec:limitations}). The protocol is the part that can carry over to experiments a team runs itself; the numbers here cannot.

\paragraph{RQ1: can a team predict a new experiment's result from its concluded experiments?} Partly. Tests from the past give a better prediction of how the next result will be distributed than a reference that draws on no past effect. A heavy-tailed base rate does that almost as well as the fitted prior, and a lead of the prior over it is not shown. On which test will be significant, the prior's skill over the base rate is under one percent in pooled form, and it is not shown in any single cell.
\begin{itemize}
% src: $RUN6/by-origin-upworthy-20261106-2dc461f52ea4.json pooled difference 0.0069 (0.0030, 0.0106), origin_level difference 0.0036 (-0.0120, 0.0191), origins 7, by_origin_above_zero 3, by_origin_below_zero 1, leave_one_origin_out[left_out=2014-07-01] difference 0.0008 (-0.0044, 0.0060), leave_one_origin_out_above_zero 6 of 7 (tables/fmt-by-origin.tex; tables/glance.tex); $RUN6/ablations/backtest-upworthy-20261106-2dc461f52ea4--exposure-planned+base-density-normal.json tables.aggregate[cell=ALL,rung=P0].logscore_vs_p1 -0.1197 (-0.1357, -0.1061), negated: P1 minus P0 in the first specification 0.120; $RUN6/ablations/backtest-upworthy-20261106-2dc461f52ea4--base-density-normal.json tables.aggregate[cell=ALL,rung=P0].logscore_vs_p1 -0.1070 (-0.1232, -0.0920), negated 0.107 (exposure fixed, normal base rate); $RUN6/ablations/backtest-upworthy-20261106-2dc461f52ea4--exposure-planned.json tables.aggregate[cell=ALL,rung=P0].logscore_vs_p1 -0.0196 (-0.0232, -0.0155), negated 0.020 (realized traffic, Student-t base rate) (tables/fmt-base-density.tex)
\item The main run fixes each row's exposure at a value known at the origin and scores the base rate with a Student-t density. In it the prior's pooled lead on the log score is 0.007 nats, and its interval by test lies above zero (Table~\ref{tab:glance}). That interval resamples tests while the seven scored quarters stay held fixed, so what it describes is these rows. With the quarters taken as the units, the interval contains zero, so a lead of the prior over the base rate is not shown; seven is a small count of values, and that interval treats them as exchangeable (Table~\ref{tab:fmt-by-origin}). Read origin by origin, the paired interval lies above zero at three origins and below zero at one, and the pooled lead depends on July 2014: leave that origin out and its interval spans zero. The first specification read each row's own realized traffic, a figure known only after launch, and scored the base rate with a normal density, and in it the prior led by 0.120 nats. That lead came from the normal density and from the realized traffic: with the normal density alone the lead is 0.107 nats, and with realized traffic alone it is 0.020 nats (Table~\ref{tab:fmt-base-density}).
% src: $RUN6/no-history-upworthy-20261106-2dc461f52ea4.json P0.minus_reference difference 0.7504 (0.6388, 0.8782), P0.by_origin_above_zero 7, P1.minus_reference difference 0.7573 (0.6483, 0.8840), P1.by_origin_above_zero 7 (tables/fmt-no-history.tex; the reference is Normal(0, s_star^2) and reads no ledger); $RUN6/learning-curve-upworthy-20261106-2dc461f52ea4.json rows[fraction=0.01,rung=P1] n_history_tests 26; paired_vs_full[fraction=0.01,rung=P1] logscore_minus_full -0.1123 (-0.1405, -0.0845); paired_vs_full[fraction=0.01,rung=P0] logscore_minus_full -0.1073 (-0.1390, -0.0773); final_origin 2015-01-01 (one origin); tables/fmt-learning-vs-full.tex; $RUN6/learning-curve-origins/learning-curve-origins-upworthy-20261106-2dc461f52ea4.json origins 7, worse_at_quarter 7
\item The fitting history is the same for both forecasters, so the difference above does not capture the contribution of history. A single comparison does capture it, set against a floor: a reference built on no past effect at all, which also assumes that no test carries an effect. At each origin, both the base rate and the prior come out ahead of it, and the base rate by itself already contains almost the whole gain (Table~\ref{tab:fmt-no-history}). Because the reference treats tests as if none of them differed, it does not stand as a serious opponent, and no score was computed for a fixed heavy-tailed density picked without reference to the archive. What more history adds is measured separately, first at a single origin. There, P1 fitted on roughly 26 tests scores 0.112 nats per row under P1 fitted on the full history, and the interval stays clear of zero (Table~\ref{tab:fmt-learning-vs-full}). The base rate gives up 0.107 nats on those same subsamples, so the gain from history is not something the prior alone produces. Run again at every scored origin, the comparison keeps its sign: with a quarter of the history the prior scores worse at all seven origins (Table~\ref{tab:fmt-learning-origins}).
\item Its skill at saying which test will be significant is under one percent of the base rate's score, and it comes from the level of the forecast and not from telling tests apart.
% src: $RUN6/backtest-upworthy-20261106-2dc461f52ea4.json tables.aggregate[cell=ALL,rung=P1] brier_skill 0.0069 (0.0041, 0.0100), adjusted_lo 0.0013, adjusted_hi 0.0125, family_size 52; tables.aggregate[rung=P1] adjusted interval by cell: headline_only (-0.0005, 0.0115), image_only (-0.0010, 0.0221), both (-0.0184, 0.0132) (tables/fmt-p1-cells.tex); tables.score_by_year[cell=ALL,rung=P1] adjusted (-0.0067, 0.0707), (-0.0018, 0.0089), (-0.0022, 0.0264) (tables/fmt-drift-year.tex); $RUN6/paper-extras-upworthy-20261106-2dc461f52ea4.json exposure_realized.p_sig_constant_within_origin_and_cell.fixed.P1 strata_with_one_value 21, origin_cell_strata 21; exposure_realized.auc_event_within_ties_exact.fixed.P1 0.5 (tables/fmt-exposure-fixed.tex)
\item When the event is set at realized precision, which is the main reading, the prior's pooled Brier skill is 0.0069, and its family-adjusted interval lies above zero. That skill is not shown in any single cell or in any single year, where each adjusted interval spans zero (Tables~\ref{tab:fmt-p1-cells} and~\ref{tab:fmt-drift-year}). Within an origin and cell the prior gives every test one probability, as the base rate does, so there it cannot tell one test from another. The pooled skill therefore reflects a better level for each cell and quarter (Table~\ref{tab:fmt-exposure-fixed}).
% src: $RUN6/paper-extras-upworthy-20261106-2dc461f52ea4.json gates_by_stratum.strata.ALL.own_gate.n_400 300 (the gate for a skill of 0.10; src/mootcourt/card.py prints n_400 as the gate for a skill of 0.10 and n_200 for 0.20); gates_by_stratum.strata.ALL.dependence.n_tests 2949, rungs.P1.verdict "decidable at skill 0.20 and 0.10", rungs.P1.brier_skill 0.0069 (below 0.01)
For a skill of 0.10, the archive clears the gate. But that gate assumes calibration, which the prior fails, and it is sized for skills far above the ones measured here, which are under 0.01 (Section~\ref{sec:gates-results}). Clearing the gate therefore certifies nothing about a skill of the measured size.
% src: $RUN6/backtest-upworthy-20261106-2dc461f52ea4.json tables.aggregate[cell=ALL,rung=P1] brier_skill_planned 0.0127 (0.0096, 0.0155); $RUN6/include-randomization-window/backtest-upworthy-20261114-4bfb354b87e3.json tables.aggregate[cell=ALL,rung=P1] brier_skill_planned 0.0103 (0.0075, 0.0130) (tables/fmt-window.tex); $RUN6/ablations/backtest-upworthy-20261106-2dc461f52ea4--base-recency-quarter.json tables.aggregate[cell=ALL,rung=P1] brier_skill_planned 0.0072 (0.0048, 0.0093); $RUN6/ablations/backtest-upworthy-20261106-2dc461f52ea4--base-recency-cell-quarter.json the same key 0.0076 (0.0047, 0.0107) (tables/fmt-base-recency.tex)
\item Setting the event at planned precision gives the prior a pooled skill of 0.0127, with its unadjusted interval above zero. The interval stays above zero when the flagged randomization window is brought in (Table~\ref{tab:fmt-window}; Section~\ref{sec:limitations}). It also stays above zero when the base rate is handed a quarter term of its own, where the skill falls to 0.0072 (Table~\ref{tab:fmt-base-recency}).
% src: $RUN6/paper-extras-upworthy-20261106-2dc461f52ea4.json exposure_realized.auc_event_within_ties_exact realized.P1 0.525, fixed.P1 0.5; exposure_realized.by_cell.ALL.realized.P1 brier_skill 0.0023 (-0.0012, 0.0061) (tables/fmt-exposure-fixed.tex; the report $RUN6/ablations/backtest-upworthy-20261106-2dc461f52ea4--exposure-planned.json)
\item Both readings change when each forecast reads its own row's realized traffic, which is a diagnostic because no forecast issued before launch could have that figure. The prior then ranks tests slightly within an origin and cell, with an AUC of 0.525 against 0.500 in the main run. Its pooled Brier skill is then 0.0023, and the interval spans zero (Table~\ref{tab:fmt-exposure-fixed}). Why the skill is higher with exposure fixed was not tested, and Table~\ref{tab:fmt-assurance} may explain it.
% src: $RUN6/paper-extras-upworthy-20261106-2dc461f52ea4.json drift_matched_null.rungs.P1.pooled current_null.percentile 97.3, mean_forecast 0.094, event_rate 0.085; exposure_realized.p_sig_constant_within_origin_and_cell.fixed.P1 strata_with_one_value 21 of 21 (one probability per origin and cell)
\item In the form in which they were issued, its forecasts fail the calibration test: the error lands at percentile 97.3 of the clustered null described in Section~\ref{sec:protocol}, where the outcomes belonging to a single test are drawn as one. The failure is mostly one of level. The mean forecast is 0.094 against a realized rate of 0.085, and the prior gives one probability per origin and cell, so what it misses is the rate of the coming quarter in each cell (Section~\ref{sec:calibration-results}).
\item Nobody knows the answer for the archive of product experiments, because that archive is undecidable at every gate.
\end{itemize}

\paragraph{RQ2: how, and how well?} Most of the log-score gain measured here comes from a heavy-tailed density fitted cell by cell to results from earlier on, while a prior fitted with measurement error contributes little on top, and even that contribution is not shown when quarters serve as the units. The method replays in time order and reads each interval. Beyond the prior, the evidence runs as follows.
\begin{itemize}
% src: $RUN6/backtest-upworthy-20261106-2dc461f52ea4.json tables.aggregate[cell=ALL,rung=P2] logscore_vs_p1 -0.0002 (-0.0008, 0.0002), brier_skill 0.0081 (0.0051, 0.0114), adjusted (0.0021, 0.0141); tables.aggregate[cell=ALL,rung=P1] brier_skill 0.0069, adjusted (0.0013, 0.0125); $RUN6/paper-extras-upworthy-20261106-2dc461f52ea4.json guard_pairs.summary pairs_admitted 4, pairs_in_selection 21
\item On the pooled log score, precedents sit level alongside the prior, with the paired difference spanning zero (Section~\ref{sec:ladder-results}). For Brier skill, the adjusted intervals of P1 and of P2 measured against the base rate both lie above zero, and the main backtest contains no paired interval between the two of them. For most origin and cell pairs, the guard admits no weight and the rung is simply the prior, so a share of the tie comes from the construction itself. Their ranking of tests runs slightly ahead of the prior's.
% src: $RUN6/paper-extras-upworthy-20261106-2dc461f52ea4.json ranking_paired.budgets["0.25"].differences.P2_minus_P1 estimate 0.0206 (0.021), adjusted_lo 0.0102, adjusted_hi 0.0310; budgets "0.10" adjusted (0.0102, 0.0280) and "0.50" adjusted (0.0065, 0.0226) (tables/fmt-ranking-paired.tex)
When the pairing is by test, the precedent rung captures an extra 0.021 of the available lift beyond the prior at a quarter of the tests, and the adjusted intervals sit clear of zero at a tenth, a quarter and a half (Table~\ref{tab:fmt-ranking-paired}).
% src: $RUN6/paper-extras-upworthy-20261106-2dc461f52ea4.json ranking_paired.budgets["0.25"].differences.P2_minus_arm_count -0.0397 (-0.0596, -0.0192), budgets["0.50"].differences.P2_minus_arm_count -0.0573 (-0.0794, -0.0361), unadjusted (adjusted_comparisons lists only the two differences from P1); Section 6.5
The edge may come from the number of arms and not from precedents: P2 captures less than a rule that counts only arms, both at a quarter and at a half of the tests, and each of those intervals lies below zero (Section~\ref{sec:decision-results}).
% src: $RUN6/paper-extras-synthetic-20261103-bf4fd286e2b2.json positive_control_rungs.runs[n_tests=11100].by_cell["ALL"].rungs[rung=P2].recovered_fraction_logscore 0.09, by_cell["headline_only"].rungs[rung=P2].recovered_fraction_logscore 0.03 (inputs.seed 0, the first ledger of 11,100 tests; $RUN6/positive-control-seeds.json seeds 5, min.recovered_pooled 0.0758, max.recovered_pooled 0.0919); by_cell["ALL"].oracle gain_logscore_vs_p1 0.0197 (0.0177, 0.0219), which the backtest reads with an interval clear of zero
\item The precedent rung manages to recover only a small share of a planted signal. Under a positive control run with the same rolling design, scoring picks up a model that was told the signal, while P2 gets back 0.09 of that model's log-score gain in pooled form and 0.03 within the headline-only cell. Both figures come from the first simulated ledger. Over five ledgers that differ in their seed, the pooled share stays between 0.08 and 0.09 (Table~\ref{tab:fmt-control-seeds}). So an interval sitting around zero does not refute a precedent signal of that size, nor does it show that the rung could use one (Section~\ref{sec:synthetic-results}).
% src: $RUN6/backtest-upworthy-20261106-2dc461f52ea4.json tables.decision_value[origin=ALL,budget=0.25,policy=p_win,rung=text_lr] diff_matched 0.046, diff_matched_adjusted_lo 0.019, diff_matched_adjusted_hi 0.073; tables.decision_value[origin=ALL,budget=0.25,policy=arm_count] captured_fraction 0.3093 against policy=p_win rung=text_lr 0.3070; $RUN6/paper-extras-upworthy-20261106-2dc461f52ea4.json ranking_paired.budgets["0.25"].differences.text_lr_minus_arm_count -0.0023 (-0.0268, 0.0221), budgets["0.10"] (-0.0079, 0.0247), budgets["0.50"] (-0.0474, 0.0031); differences.P2_minus_arm_count budgets["0.10"] (-0.0259, 0.0009), budgets["0.25"] (-0.0596, -0.0192), budgets["0.50"] (-0.0794, -0.0361) (tables/fmt-ranking-paired.tex)
\item Among the models, the text classifier comes out on top in ranking. At a quarter of the tests it beats the arm-count-matched random order by 0.046 of captured lift, and its adjusted interval stays clear of zero. No tuning was done on it, and a single split produced the number. Set against a rule whose ranking uses arm count alone and never looks at any forecast, its paired difference at that budget is minus 0.002, and the interval spans zero under each of the three budgets (Table~\ref{tab:fmt-ranking-paired}). So no model is shown to rank tests better than counting arms.
% src: $RUN6/backtest-upworthy-20261106-2dc461f52ea4.json tables.aggregate[cell=ALL,rung=text_lr] brier_skill 0.0070 (-0.0005, 0.0140), brier_skill_planned 0.0153 (0.0069, 0.0238); $RUN6/include-randomization-window/backtest-upworthy-20261114-4bfb354b87e3.json tables.aggregate[cell=ALL,rung=text_lr] brier_skill_planned 0.0047 (-0.0018, 0.0116) (tables/fmt-window.tex); header.challengers empty in the two base-recency reports
\item Skill on significance for the text classifier turns on how the event is defined: at realized precision the interval spans zero, while at planned precision in the main run it sits above zero, unadjusted. Once the randomization window is brought in, that interval spans zero as well (Table~\ref{tab:fmt-window}), and no scoring was done of the classifier against a base rate carrying a quarter term.
% src: $RUN6/paper-extras-synthetic-20261103-bf4fd286e2b2.json positive_control_decision_value (tables/fmt-ranking-paired.tex: planted ledgers, text classifier minus P1 -0.0119 and -0.0199)
\item On skill the text classifier fails the planted control, and it also ranks under P1 there, so the control neither validates nor explains its value on the real archive, and the question of what it is reading there stays open.
\item On this archive abstention does not improve what gets issued. Within an origin and cell the abstained rows score above the issued rows on Brier skill, and the interval for that difference excludes zero (Table~\ref{tab:fmt-abstention-composition}). So abstention remains a label, one that says which trigger fired while making no claim about how good the forecast is (Section~\ref{sec:abstention-results}). Next to each forecast the card prints the most recent gated skill for its cell, or, when the cell is undecidable, the pooled figure with its label.
% src: $RUN6/paper-extras-upworthy-20261106-2dc461f52ea4.json sizing_evaluation.test_and_roll.by_audience["20"].paired_differences.prior_minus_base_rate estimate -0.0016, adjusted_lo -0.0042, adjusted_hi 0.0010 (tables/fmt-sizing.tex)
\item The sizing evaluation ran over a synthetic horizon and involved no money (Table~\ref{tab:fmt-sizing}).
% src: $RUN6/paper-extras-upworthy-20261106-2dc461f52ea4.json sizing_evaluation.test_and_roll primary_multiple 20, by_audience["20"].paired_differences.prior_minus_base_rate estimate -0.0016, adjusted_lo -0.0042, adjusted_hi 0.0010; prior_minus_power adjusted (-1.349, -1.083), base_rate_minus_power -1.214 (-1.314, -1.102), relative regret, lower is better; tables/fmt-sizing.tex fixed power rule 1.568 at the primary multiple, equal to test the whole audience, against never test 1.000
At the primary horizon, where the audience runs to 20 times the size of the median two-arm test in the cell, the size from the prior shows no measurable edge over the size from the base rate. Each of them beats a rule that was sized for a fixed power, but the win comes only because the rule reaches its cap at that horizon and ends up testing the entire audience, an outcome that loses more than not testing at all (Section~\ref{sec:sizing-results}). When the comparison is with the base rate, the sign flips as the horizon changes and as the assumed world changes. No claim is made in this paper that taking the size from the forecast does better than a base rate.
\end{itemize}

%% tmlr:move-begin production
%% tmlr:note production What improved the prediction (RQ3) is read in Appendix~\ref{sec:production}, with notes on using the package in production, where it has not been deployed. Of the things put to the test, the biggest gain, on both forecasters, comes from a heavy-tailed density, and the bulk of the settings do not matter.
%% tmlr:title production Further discussion: what improves the prediction, and production use
\label{sec:production}
\paragraph{RQ3: what improves the prediction?} Of the things put to the test, a handful of changes make a difference and the bulk of the settings do not matter.
\begin{itemize}
% src: $RUN6/ablations/backtest-upworthy-20261106-2dc461f52ea4--prior-normal.json tables.aggregate[cell=ALL,rung=P0] logscore_vs_p1 0.0700 (0.0502, 0.0936): a normal prior is 0.070 nats below the Student-t base rate (tables/fmt-prior-settings.tex); $RUN6/ablations/backtest-upworthy-20261106-2dc461f52ea4--base-density-normal.json tables.aggregate[cell=ALL,rung=P0] logscore_vs_p1 -0.1070 (-0.1232, -0.0920): a normal base rate is 0.107 nats below the Student-t prior (tables/fmt-base-density.tex)
\item The biggest gain, on both forecasters, comes from a heavy-tailed density. A prior with a normal density scores 0.070 nats below the heavy-tailed base rate, and a base rate with a normal density scores 0.107 nats below the heavy-tailed prior (Tables~\ref{tab:fmt-prior-settings} and~\ref{tab:fmt-base-density}). So what this states concerns the shape of results at this scale, and it says nothing about how the prior is designed. One in-sample diagnostic, reading a rate that only becomes known after launch, points to much of what the Student-t picks up being a mixture across the click-through level. That diagnostic leaves the gain unsplit, and no claim is made here that the effects are themselves heavy-tailed (Table~\ref{tab:fmt-prior-family}).
\item For the reasons set out under RQ2, the ranking of the untuned text classifier on this single split counts as a lead, not as a change that has been shown to make prediction better.
\item Calibration would get better if a forecast could anticipate the level. Section~\ref{sec:calibration-results} shows how large that error is. The shift it fits uses the realized rate, so no method in this paper delivers the gain.
% src: $RUN6/ablations/backtest-upworthy-20261106-2dc461f52ea4--lambda-1.0.json tables.aggregate[cell=ALL,rung=P2] logscore_vs_p1 -0.01272
\item Setting a precedent weight to one, and holding it there, worsens the score by 0.013 nats when it is set against the prior.
% src: $RUN6/ablations/backtest-upworthy-20261106-2dc461f52ea4--p2-guard-off.json tables.aggregate[cell=ALL,rung=P2] logscore_vs_p1 -0.00314 (-0.00441, -0.00186); $RUN6/ablations/backtest-upworthy-20261106-2dc461f52ea4--p2-selection-loto+p2-guard-off.json tables.aggregate[cell=ALL,rung=P2] logscore_vs_p1 0.00026 (-0.00098, 0.00162)
With the guard off, the selected weight still carries a cost of 0.003 nats under the default selection rule, and turning the guard on takes that cost away. When selection is leave-one-test-out, the paired difference of P2 from the prior spans zero while the guard is off, so the need for the guard attaches to the default rule and not to every rule (Section~\ref{sec:ablations}).
% src: $RUN6/ablations/backtest-upworthy-20261106-2dc461f52ea4--features-text.json tables.aggregate[cell=ALL,rung=P2] logscore_vs_p1 -0.00148 (-0.00223, -0.00081); $RUN6/ablations/backtest-upworthy-20261106-2dc461f52ea4--kappa-0.0.json tables.aggregate[cell=ALL,rung=P1] logscore -0.7733 and $RUN6/ablations/backtest-upworthy-20261106-2dc461f52ea4--kappa-80.0.json -0.7675, the two ends of the grid (tables/fmt-prior-settings.tex prints -0.773 and -0.768; the report holds no paired interval between settings of kappa); kappa-0.0 tables.aggregate[cell=ALL,rung=P0] logscore_vs_p1 -0.0028 (-0.0075, 0.0023), the lead of the prior with no pooling, which spans zero
\item Bandwidth, recency and the selection rule barely shift P2's score, and using text features on their own puts it under the prior, with an interval that stays clear of zero (Table~\ref{tab:fmt-p2-settings}). Across its grid, the pooling strength takes the prior's log score from minus 0.773 up to minus 0.768. That shift is nearly as large as the prior's whole lead over the base rate, and with no pooling the interval for that lead spans zero; no paired interval compares the settings with each other (Table~\ref{tab:fmt-prior-settings}).
% src: $RUN6/learning-curve-upworthy-20261106-2dc461f52ea4.json rows rung=P1 skill 0.013 at n_fit_rows 2369.2 and at 9447 (tables/fmt-learning.tex); paired_vs_full[fraction=0.25,rung=P1] brier_gain_over_full -0.000003 (-0.000060, 0.000064), logscore_minus_full -0.0067 (-0.0093, -0.0043); paired_vs_full[fraction=0.5,rung=P1] brier_gain_over_full -0.000023 (-0.000064, 0.000026), logscore_minus_full -0.0013 (-0.0025, -0.0003), logscore_seed_sd 0.0029, above the mean gap, and the interval holds the subsamples fixed (tables/fmt-learning-vs-full.tex); $RUN6/learning-curve-origins/learning-curve-origins-upworthy-20261106-2dc461f52ea4.json rows quarter brier_hi below zero at the five origins from 2013-04-01 to 2014-10-01
\item At the origin of the first curve, measured against the full history, the paired Brier difference for the prior spans zero once a quarter of the history is in, a result that does not amount to showing the gain ends there (at the five earliest origins that difference lies below zero, Table~\ref{tab:fmt-learning-origins}), and the mean gap on the log score at half of it is small and not established (Tables~\ref{tab:fmt-learning} and~\ref{tab:fmt-learning-vs-full}). Scoring of the skill curve uses the base rate from the full history, and that rate is stale at that origin, so the curve does not show the amount of skill the prior carries (Section~\ref{sec:ablations}).
\item Not tested, so these are candidates and not findings: recording planned exposure, stronger text models, choosing the width of the precedent rung's forecast on held-out rows, and team forecasts frozen before launch.
\end{itemize}

\paragraph{Using this in production.} No production experimentation platform has ever run this system in deployment, and a replay is what the evidence rests on. A team gets from it a means of putting forecasts to the test against the team's own ledger. Aggregates are all it requires, and at no point does the ledger go outside the organization. The order in which it is used is fixed.
\begin{enumerate}
% src: $RUN6/paper-extras-upworthy-20261106-2dc461f52ea4.json gates_by_stratum.strata.ALL rungs.P1.verdict "decidable at skill 0.20 and 0.10" (the two skills the gates are sized for); $RUN6/backtest-upworthy-20261106-2dc461f52ea4.json tables.aggregate[cell=ALL,rung=P1] brier_skill 0.0069 (0.0041, 0.0100)
\item Bring in the experiments that have concluded, and run the backtest over those before anything else. Where forecasts fail on that ledger, and the number of tests that a claim of a stated skill requires under the assumptions of the gate, come from that report and not from the numbers in this paper. Certifying a forecast is something it has not been shown to do: the sizing of its gates targets skills of 0.10 and 0.20 with calibration assumed, and the skills measured in this work sit far below those.
% src: $RUN6/backtest-upworthy-20261106-2dc461f52ea4.json tables.decision_value[origin=ALL,policy=p_win,rung=P2] diff_matched_adjusted_lo 0.0009 at budget 0.10, -0.0088 at 0.25, -0.0175 at 0.50; $RUN6/paper-extras-upworthy-20261106-2dc461f52ea4.json ranking_paired.budgets differences.P2_minus_arm_count (-0.0259, 0.0009) at "0.10", (-0.0596, -0.0192) at "0.25", (-0.0794, -0.0361) at "0.50"
\item For a claim that is decidable, the card sets out ahead of launch the result that the ledger points one toward expecting. The probabilities issued on this archive failed the calibration test, so before a team acts on the probability on a card it should look at the team's own reliability table. When candidate tests are ranked by the card, the ranking beats a random order matched on arm count at one tenth of the tests and nowhere else, after adjustment. One rule, counting arms and nothing else, captured lift beyond what the card captured, and its paired interval sits below zero at a quarter of the tests and at a half, so a team ought to check the ranking against the team's own backtest before putting any reliance on it. A card can also suggest a size, and any value it has over a base rate is not shown.
\item Outcomes are to be joined weekly. When support is thin, or novelty is present, or the reliability record has gone stale or drifted, or the decision is unstable, the card abstains, and that abstention labels the evidence without being shown to make what is issued any better.
% src: $RUN6/replay-upworthy-20261106-2dc461f52ea4.json totals.n_flagged_weeks 6, header.n_weeks 111, weeks[].drifted_cells both only (tables/replay-upworthy.tex); $RUN6/backtest-upworthy-20261106-2dc461f52ea4.json tables.aggregate[cell=both,rung=P1] n_tests 173, the fewest of the three cells
Drift was flagged by the weekly loop in 6 of 111 weeks on this archive, and every one of those weeks fell in the smallest cell (Section~\ref{sec:drift-results}). The rule sets a cell's log score against the base rate's, so a change missed jointly by both of them may escape it, such as the year-to-year fall in the significant-positive rate here. Alongside the flag, a team should keep an eye on the event rate in each cell.
Anything above the prior must earn its place on a paired, gated comparison.
\end{enumerate}
Four limits apply to that use.
\begin{itemize}
\item Randomized A/B and A/B/n tests are what the evidence covers, and it comes mainly from a single publisher, so any other design has to be treated by a team as untested (Section~\ref{sec:limitations}, Scope).
\item Weak evidence and leakage in a backtest are what the protocol turns up. Defects in a change are not what it finds: bugs, instrumentation errors and harm stay the job of tests, guardrails and the live experiment.
\item No promise of improvement on product experiments is made in this paper, because the single archive of them is undecidable.
\item A team ranking by the forecast ends up selecting the ledger it will have later, so the rows should note whether a card had an influence on the test (Section~\ref{sec:limitations}).
\end{itemize}
%% tmlr:move-end production
